\section{Complexity of Nonconvex Constrained DC Optimization}
\label{sec:complexity-analysis-NCDCA}
In this section, we analyze the complexity of the proposed algorithmic framework for solving the constrained DC problem \eqref{prob:main}. Beyond the assumptions stated in \eqref{prob:main}, we impose the following additional assumptions throughout this section:

\subsection{Assumptions}
Given nonconvex constraints, we impose the following nonsingularity condition on the constraints to ensure the feasibility of solutions. Constraint qualifications are often required in solving constrained optimization problems; see~\cite{shi2025momentum,lu2024variancereducedfirstordermethodsdeterministically,liu2025single,curtis2024worst}.
\begin{assumption}
    \label{Ass: constraints conditions} 
    {For the iterate sequence $\{\bx^k\}_{k\in\NN}$ generated by the algorithms,} there exists a constant $\gamma_c > 0$ such that
    \begin{align*}
        \|\nabla \bc_{\calE}(\bx^k)\bc_{\calE}(\bx^k)
        +\nabla \bc_{\calI}(\bx^k)[\bc_{\calI}(\bx^k)]_+\|
        \geq
        \gamma_c
        \left(\|\bc_{\calE}(\bx^k)\|^2+\|[\bc_{\calI}(\bx^k)]_+\|^2\right)^{1/2},
        \quad \forall\, k \geq 0.
    \end{align*}
\end{assumption}

\subsection{Convergence Criterion}
\label{subsec:Convergence Criterion of KKT}

The goal of this subsection is to establish the certificate generated by
the algorithmic framework below. For each iteration, we define
the transported same-point tolerances
\[
    \delta_\phi^{k+1}:=\epsilon_\phi^{k+1},
\]
\[
    \delta_g^{k+1}
    :=
    \epsilon_g^k
    +g(\bx^{k+1})-g(\bx^k)
    -\langle \bb^k,\bx^{k+1}-\bx^k\rangle,
\]
and the certificate residual
\[
    R_{k+1}
    :=
    \frac{1}{\gamma_k}\|\bx^{k+1}-\bx^k\|
    +\alpha_{k+1}.
\]
The main result of this subsection is the bound
\[
    \operatorname{dist}
    \left(
        \partial_{\delta_\phi^{k+1}}\phi(\bx^{k+1}),
        \partial_{\delta_g^{k+1}}g(\bx^{k+1})
    \right)
    \le R_{k+1}.
\]
This certificate separates three effects: the $\phi$-side oracle error,
the transported $g$-side Fenchel error, and the proximal optimality
residual.

% \subsection{Deterministic Case}
% \label{subsec:deterministic-case}

% In the deterministic setting, we instantiate the framework with exact
% component certificates and a controlled proximal optimality residual.  In
% the cleanest case,
% \[
%     \epsilon_\phi^{k+1}=0,\qquad
%     \epsilon_g^k=0,\qquad
%     \alpha_{k+1}=0.
% \]
% The analysis then shows that a suitable output iterate satisfies
% \[
%     \delta_\phi^{\hat k+1}=0,\qquad
%     \delta_g^{\hat k+1}=O(\epsilon),
%     \qquad
%     R_{\hat k+1}=O(\epsilon)
% \]
% within
% \[
%     K=O(\epsilon^{-2})
% \]
% iterations.  This case demonstrates that nonsmoothness of the subtracted
% component can be handled through the transported $g$-side Fenchel
% certificate.

% \subsection{Finite-Sum Case}
% \label{subsec:finite-sum-case}

% We next consider the finite-sum model, where the stochastic structure
% appears in a smooth component, for example
% \[
%     \phi(\bx)=f(\bx)+h(\bx),
%     \qquad
%     f(\bx)=\frac{1}{N}\sum_{i=1}^N f_i(\bx).
% \]
% Variance-reduced estimators such as PAGE, SARAH, or SPIDER affect the
% certificate through the residual term $R_{k+1}$, or equivalently through
% the mismatch between the approximate smooth gradient and the true
% gradient.  The $g$-side tolerance remains governed by the transported
% Fenchel error
% \[
%     \delta_g^{k+1}
%     =
%     \epsilon_g^k
%     +g(\bx^{k+1})-g(\bx^k)
%     -\langle \bb^k,\bx^{k+1}-\bx^k\rangle.
% \]
% This subsection clarifies how finite-sum variance reduction can be
% combined with the two-sided certificate, and how the resulting guarantee
% differs from one-sided DCA gap guarantees.

% \subsection{Online Case}
% \label{subsec:online-case}

% Finally, we consider the online expectation setting, especially the case
% where the subtracted component $g$ is stochastic and nonsmooth.  The
% two-sided certificate exposes a bottleneck that is hidden by one-sided
% gap functions.  The $g$-side tolerance decomposes as
% \[
%     \delta_g^{k+1}
%     =
%     \epsilon_g^k
%     +
%     \underbrace{
%     g(\bx^{k+1})-g(\bx^k)
%     -\langle \bb^k,\bx^{k+1}-\bx^k\rangle
%     }_{\text{transport error}}.
% \]
% For a generic bounded-variance stochastic nonsmooth oracle, obtaining an
% old-point certificate with $\epsilon_g^k=O(\epsilon)$ typically requires
% $O(\epsilon^{-2})$ samples, while controlling the transport error through
% the step size requires $O(\epsilon^{-2})$ outer iterations.  This leads to
% the balanced online scaling
% \[
%     O(\epsilon^{-4})
% \]
% for the two-sided Fenchel certificate.  Therefore, an
% $O(\epsilon^{-3})$ online rate would require additional structure, such
% as a smooth or variance-reduced Fenchel oracle for the $g$-side.

